\begin{lemma}[Endpoint exclusion in an elementary augmented tree]
Let $T,U$ be exact tree sequences of a finite loopless multigraph, under
the partial edge coloring $\varphi$, and let $P_i$ be $(\alpha,\beta)$ exit paths of
$T$ from $a_i$ to $b_i$ for $i=1,2,3$.  Suppose $a_1\ne a_2$,
all $b_i$ belong to $U$, $U$ is elementary, $\beta$ is missing at
$b_3$, and $U$ is singleton-$\beta$-closed.  Suppose also that whenever
$\alpha$ is missing on $U$, $U$ is singleton-$\alpha$-closed.
Then none of $a_1,a_2,a_3$ belongs to $U$.
\end{lemma}
\begin{proof}
By \noderef{TreeExitPath}, each $a_i$ is outside $T$, each $b_i$ is in
$T$, and each $a_i$ misses at least one of the two colors.  In particular
$a_i\ne b_3$.  If $a_i$ belongs to $U$, elementarity
(\noderef{ColoredVertexSetElementary}) and the missing color $\beta$ at
$b_3$ prohibit $\beta$ from being missing at $a_i$.  Consequently
$\alpha$ is missing at $a_i$, hence on $U$ by
\noderef{TreeMissingColors}.  The assumed implication now makes both
colors closed on $U$.

We show that with both colors closed any two-color walk whose final
vertex belongs to $U$ has its initial vertex in $U$.  Traverse the walk
backwards.  For vertex-nodes use membership in $U$ as the invariant; for
an edge-node use the condition that every endpoint of that exact edge
belongs to $U$.  A transition from an edge-node to a vertex-node
preserves the invariant by \noderef{TwoColorIncidenceGraph}.  For a
transition from a vertex $v$ in $U$ to an edge $e$, the same definition
gives incidence and one of the two colors.  If an endpoint $w$ of $e$
were outside $U$, then $v\ne w$ and the unordered endpoint pair is
$\{v,w\}$.  Thus $e$ is a boundary edge by \noderef{BoundaryEdges},
and a boundary edge of one of the two colors by
\noderef{BoundaryColorEdges}.  This contradicts
\noderef{TreeColorSetClosed}.  Induction over the finite walk proves
the claimed invariant.

Apply this to $P_1$ and $P_2$, whose final vertices are in $U$.  Both
$a_1,a_2$ are now in $U$, and, by the first paragraph, both miss
$\alpha$.  Since they are distinct, this contradicts elementarity.
Thus the assumption that any one of the three $a_i$ is in $U$ leads
to a contradiction, proving all three exclusions.
\end{proof}
